\chapter{Extending the framework}
\label{ch:development}

This chapter summarises the rules for adding code to ATTPCROOT. The reference pages in \code{docs/contributing/}
and \code{docs/tooling/} hold the same rules in more detail.

\section{Libraries and dictionaries}

Each module is a CMake shared library with a ROOT dictionary. Its \code{CMakeLists.txt} lists the sources, the
dependencies and the LinkDef file, and registers the library and its tests:
\begin{verbatim}
set(LIBRARY_NAME MyModule)
set(SRCS MyClass.cxx MyTask.cxx)
set(DEPENDENCIES ATTPCROOT::AtData ROOT::Core FairRoot::Base)
set(TEST_SRCS MyClassTest.cxx)
attpcroot_generate_tests(${LIBRARY_NAME}Tests
   SRCS ${TEST_SRCS} DEPS ${LIBRARY_NAME})
generate_target_and_root_library(${LIBRARY_NAME}
   LINKDEF ${LIBRARY_NAME}LinkDef.h SRCS ${SRCS} DEPS_PUBLIC ${DEPENDENCIES})
\end{verbatim}
A new module directory is added to the top-level \code{CMakeLists.txt} with \code{add\_subdirectory}, and
\code{ATTPCROOT::MyModule} to the dependencies of every module that uses it.

\paragraph{LinkDef suffixes.} Every class in a \code{*LinkDef.h} must carry the right suffix:
\begin{itemize}
\item \code{ClassName +;} generates a full I/O streamer. Use it only for classes written to a ROOT file (stored in a
  branch): \code{AtHit}, \code{AtEvent}, \code{AtRawEvent}, \code{AtMCPoint}, the PID result structures.
\item \code{ClassName -!;} gives reflection only, with no streamer. Use it for everything else: tasks, algorithms,
  models, samplers. Default to \code{-!} unless the class is persisted.
\end{itemize}
A persisted class that changes its layout must increment its \code{ClassDef} version, so that older files are read
with schema evolution (for example \code{AtPIDEvent} went from version 2 to 3 when the combined PID was added; older
files read back with the new vector empty).

\paragraph{Types.} Non-persisted classes use \code{bool}, \code{int}, \code{double}, \code{std::string}; persisted
data classes use the ROOT typedefs (\code{Bool\_t}, \code{Int\_t}, \code{Double\_t}). Persisted classes in this tree
already use STL containers and \code{std::unique\_ptr} where dictionaries support them (\code{AtEvent} stores its hits
in \code{std::vector<std::unique\_ptr<AtHit>>}); follow the pattern of the owning module.

\section{Tasks}

Most processing steps are \code{FairTask}s with the same structure: \code{Init()} retrieves the input branch from the
\code{FairRootManager} and registers the output branch; \code{SetParContainers()} requests run-time parameter
containers when needed; \code{Exec()} clears the output, reads the current input and fills the output. Objects read
from the \code{FairRootManager} are owned by the framework: never delete them or wrap them in smart pointers. Macros
create tasks with \code{new} because \code{AddTask} takes ownership; inside a task, owned algorithms are usually held
in a \code{std::unique\_ptr}. A task that reads or writes documented branches must update
\code{docs/reference/branch-io-contracts.md}.

A new option on an existing task should default to the old behaviour, so that every existing production is
unchanged until someone switches it on. The PID options of Chapter~\ref{ch:pid} follow this rule.

\section{Tests}

Unit tests use GoogleTest and live next to the code they test: the tests of \code{AtHelix3D} are in
\code{AtTools/AtPID/AtHelix3DTest.cxx}. They are registered with \code{attpcroot\_generate\_tests} and run by \code{ctest}. They must not
read external files or use the network: put small inputs inline (the pad-gain parser test reads its input from an
\code{std::istringstream}). Keep algorithms free of FairRoot where possible so they can be tested directly; the PID
combination rule is a plain function, \code{AtTools::CombinePID}, for this reason. Integration tests that need data or
geometry are ROOT macros in \code{macro/tests/}, run by \code{macro/tests/runAllTest.sh}.

\section{Style}

Libraries are named after their directory (\code{AtFoo} builds \code{libAtFoo}); classes outside a namespace carry the
\code{At} prefix; data members carry an \code{f} prefix. Code is formatted with \code{clang-format-17} and the
repository \code{.clang-format} (LLVM based, 3-space indent, 120 columns); \code{scripts/formatAll.sh} formats the
changed files. \code{.clang-tidy} configures the static analysis. Pull requests target \code{develop}, are merged
fast-forward only, and must pass \code{clang-format}, \code{clang-tidy} and the unit tests. Commit messages are in the
imperative mood and at most 72 characters in the subject line.

\section{Documentation}

Two forms of documentation live in \code{docs/}: short reference pages in Markdown (\code{docs/index.md} is the entry
point) and this manual (\code{docs/manual/}). The manual is built with
\begin{verbatim}
cd docs/manual
pdflatex attpcroot_manual.tex && pdflatex attpcroot_manual.tex
\end{verbatim}
Each chapter is a separate file in \code{docs/manual/chapters/}; \code{test\_chapter.sh <name>} compiles one chapter
alone and reports errors and overfull lines. The style rules are in \code{docs/manual/STYLE.md}: plain LaTeX, every
statement backed by the code, defaults quoted from the source. When a change alters the behaviour of a component,
update its chapter in the same commit.
